\section*{Bab \ref{ch:b3:virology} --- Virologi}

\begin{solution}{exo:b3:virology:1}
Polio: IV, RNA untai positif --- sebuah RNA polimerase bergantung-RNA.
Influenza: V, RNA untai negatif --- RNA polimerasenya sendiri, yang diusung di
dalam \omterm{def:b3:virology:virus}{virion}. HIV: VI --- transkriptase balik (dan integrase). Herpes
simpleks: I, DNA untai ganda --- inangnya pada asasnya dapat memasok
segalanya, tetapi ia membawa DNA polimerase dan timidin kinasenya sendiri.
Hepatitis B: VII --- sebuah transkriptase balik untuk menyalin RNA
pragenomnya menjadi DNA. Rotavirus: III, RNA untai ganda --- sebuah RNA
polimerase bergantung-RNA di dalam partikelnya.
\end{solution}

\begin{solution}{exo:b3:virology:2}
DNA fag itu ($^{32}$P) memasuki bakterinya dan proteinnya
($^{35}$S) tinggal di luar lalu dapat digunting lepas; fag keturunannya
mewarisi $^{32}$P itu. Jadi DNA-lah yang disuntikkan fag itu dan yang
mengarahkan pembuatan fag baru. Seandainya proteinlah yang mengusung
petunjuknya, $^{35}$S akan ditemukan di dalam lalu diteruskan ke keturunannya.
\end{solution}

\begin{solution}{exo:b3:virology:3}
$84/(0.1\times 10^{-7}) = 8.4\times 10^{9}$ satuan pembentuk plak per mL.
Mikroskopnya mencacah tiap partikel, termasuk \omterm{def:b3:virology:virus}{kapsid} kosong, partikel
bergenom cacat, dan partikel yang gagal melekat atau menyuntik; hanya
partikel yang menuntaskan sebuah infeksi yang membuat plak.
\end{solution}

\begin{solution}{exo:b3:virology:4}
Pelekatan: penghambat pemasukan (maravirok), antibodi penetral.
Pemasukan atau peleburan: penghambat peleburan, penghalang pengasaman \omterm{def:b3:membrane-traffic:endocytosis}{endosom}.
Pengupasan mantel: obat pengikat \omterm{def:b3:virology:virus}{kapsid} (plekonaril; lenakapavir).
Pengungkapan dan penggandaan: \omterm{def:b3:virology:drugs}{analog nukleosida}, penghambat transkriptase
balik dan penghambat integrase, interferon, \omterm{def:b3:rna-regulation:rnai}{interferensi RNA}. Perakitan:
\omterm{def:b3:virology:drugs}{penghambat protease} (poliprotein yang belum terpotong tidak dapat merakit).
Pelepasan: penghambat neuraminidase (influenza), dan sel~T sitotoksik yang
membunuh selnya sebelum pelepasan.
\end{solution}

\begin{solution}{exo:b3:virology:5}
Satu protein \qty{30}{kDa} berupa $30\,000/110 = 273$ residu, yaitu $818$
nukleotida: \qty{0.82}{kb}, yakni \qty{18}{\%} dari genom \qty{4.5}{kb}.
Menyandi $180$ salinannya secara terpisah akan memakan \qty{147}{kb},
yakni tiga puluh tiga kali seluruh genomnya; protein \omterm{def:b3:virology:virus}{kapsidnya} sendiri
(\qty{5.4}{MDa}) berbobot empat kali genomnya (\qty{1.5}{MDa}).
\end{solution}

\begin{solution}{exo:b3:virology:6}
$V = 10^{5}(3e^{-0.5t} - 0.5e^{-3t})/2.5$: $t = 0.5$: $8.9\times
10^{4}$; $t = 2$: $4.4\times 10^{4}$; $t = 10$: $8\times 10^{2}$. Di bawah
$50$ ketika $1.2\times 10^{5}e^{-0.5t} = 50$: $t = 2\ln(2400) \approx
\qty{16}{d}$.
\end{solution}

\begin{solution}{exo:b3:virology:7}
$uL = 0.36$; pecahan bebas galat $e^{-0.36} = 0.70$. Tiga kali lipat: $uL =
1.08$, $e^{-1.08} = 0.34$ --- kebanyakan keturunannya kini mengusung mutasi
dan urutan terbugarnya terencerkan ke arah ambangnya. Sebuah koronavirus pada
\qty{30}{kb} dengan $u = 10^{-4}$ akan punya $uL = 3$ dan hanya \qty{5}{\%}
salinan bebas galat, yakni melewati ambangnya; eksonukleasenya membawa $u$ ke
$10^{-5}$, $uL = 0.3$, yakni tiga perempat bebas galat.
\end{solution}

\begin{solution}{exo:b3:virology:8}
Hanyutan: mutasi titik pada hemaglutinin dan neuraminidase menimbun di bawah
seleksi antibodi, sehingga galur tiap musim sebagian baru dan imunitas tahun
lalu sebagian usang. Pergeseran: infeksi bersama satu sel oleh sebuah galur
manusia dan sebuah galur hewan menyusun ulang kedelapan segmennya, dan sebuah
subtipe hemaglutinin dari burung atau babi muncul di dalam sebuah \omterm{def:b3:virology:virus}{virus} yang
teradaptasi pada manusia. Tak seorang pun punya antibodi terhadap subtipe yang
baru itu, sehingga seluruh populasi rentan sekaligus --- yakni syarat sebuah
pandemi, yang terpenuhi pada tahun 1918, 1957, 1968, dan 2009 oleh sebuah
hemaglutinin baru tiap kali.
\end{solution}

\begin{solution}{exo:b3:virology:9}
Timidin kinase herpesvirus memfosforilasi asiklovir, yang hampir tidak
disentuh kinase sel; kinase sel lalu merampungkan trifosfatnya, yang
disisipkan DNA polimerase \omterm{def:b3:virology:virus}{virus} dan, karena tidak punya hidroksil 3$'$,
menamatkan rantainya. Sel yang tidak terinfeksi tidak pernah membuat bentuk
aktifnya, dan polimerase \omterm{def:b3:virology:virus}{virus} lebih menyukainya daripada polimerase sel.
Kekebalan: hilangnya atau berubahnya timidin kinase \omterm{def:b3:virology:virus}{virus} (yang
terlazim), atau mutasi pada polimerase \omterm{def:b3:virology:virus}{virus} yang menyingkirkan
analognya.
\end{solution}

\begin{solution}{exo:b3:virology:10}
Dengan $T$ tetap, persamaannya linear dalam $(I, V)$ dengan matriks
$\begin{pmatrix} -\delta & \beta T \\ p & -c \end{pmatrix}$, yang jumlah
diagonalnya $-(\delta + c) < 0$ dan determinannya $\delta c - \beta Tp$. Nilai
eigennya berjumlah negatif, sehingga salah satunya positif tepat ketika
determinannya negatif: $\beta Tp > c\delta$, yakni $R_{0} = \beta
Tp/(c\delta) > 1$. Faktornya: $p/\delta$ adalah banyaknya \omterm{def:b3:virology:virus}{virion} yang dibuat
sebuah sel terinfeksi selama hidupnya; $\beta T/c$ adalah banyaknya sel yang
diinfeksi sebuah \omterm{def:b3:virology:virus}{virion} selama hidupnya; hasil kalinya adalah banyaknya sel
terinfeksi yang dihasilkan satu sel terinfeksi.
\end{solution}

\begin{solution}{exo:b3:virology:11}
$10^{6} = 2^{20}$: dua puluh paruh hidup, $880$ bulan, sekitar $73$ tahun ---
lebih lama daripada satu kehidupan. Obat yang menghalangi penggandaan tidak
menyentuh sebuah provirus yang tidak sedang ditranskripsi; wadahnya bertahan,
dan ketika terapinya berhenti satu sel yang mengaktif kembali memulai lagi
infeksinya. Sebuah kesembuhan harus menyingkirkan atau membungkam wadah itu
selamanya: membunuh sel latennya (mengaktifkannya kembali di bawah terapi
sehingga ia mati atau dibunuh), memotong keluar atau melumpuhkan provirusnya,
atau mengganti sistem imunnya dengan yang tidak dapat diinfeksi \omterm{def:b3:virology:virus}{virus} itu.
\end{solution}

\begin{solution}{exo:b3:virology:12}
$uL = 0.03$ mutasi per genom; $10^{9}$ genom sehari mengusung $3\times
10^{7}$ mutasi di antara $9\times 10^{4}$ mutasi tunggal yang mungkin:
masing-masing muncul sekitar $300$ kali sehari. Sebuah mutan rangkap dua
tertentu muncul pada $10^{-12}$ per genom, yakni $10^{-3}$ sehari --- sekali
dalam tiga tahun; sebuah rangkap tiga pada $10^{-18}$, tidak pernah. Dua obat
dengan mutasi kekebalan yang saling bebas akan mencukupi pada asasnya, tiga
memberi selisih bagi kepatuhan yang buruk. Dua obat yang mengikat enzim yang
sama dapat dikalahkan satu mutasi yang mengubah tapaknya atau konformasi
enzimnya, dan mutasi kekebalannya tidak saling bebas; keduanya terhitung
kurang dari dua.
\end{solution}

\begin{solution}{pb:b3:virology:1}
\textbf{1.} $2\times 10^{5}\times 1.5\times 10^{4} = 3\times 10^{9}$
\omterm{def:b3:virology:virus}{virion}.
\textbf{2.} Disingkirkan $cV = 3\times 3\times 10^{9} \approx 10^{10}$ sehari;
dihasilkan sebanyak itu pula.
\textbf{3.} $I^{*} = cV^{*}/p = 9\times 10^{9}/500 = 1.8\times 10^{7}$
sel terinfeksi; yang mati $\delta I^{*} = 9\times 10^{6}$ sehari.
\textbf{4.} RNA $2\times 9700\times 330 = 6.4\times 10^{6}$ Da; protein
$2000\times 25\,000 = 5\times 10^{7}$ Da; seluruhnya $5.6\times 10^{7}$ Da
$\approx 10^{-16}$ g. Seluruh \omterm{def:b3:virology:virus}{virion} bersama: $3\times 10^{9}\times
10^{-16} = 3\times 10^{-7}$ g, yakni sepertiga mikrogram --- tiga
ribu kali lebih ringan daripada sebutir pasir.
\textbf{5.} $1.8\times 10^{7}/10^{12} = 2\times 10^{-5}$ dari kumpulannya.
Sepuluh tahun dengan $9\times 10^{6}$ kematian sehari berjumlah
$3\times 10^{10}$ sel, yakni tiga persen kumpulannya: kumpulan itu gagal
bukan karena pengurangan aritmetis melainkan karena sedasawarsa peremajaan
paksa dan pengaktifan menahun menghabiskan kemampuan menggantikannya.
\textbf{6.} $9\times 10^{6}/24 \approx 4\times 10^{5}$ bangkai yang sedang
disingkirkan setiap saat.
\textbf{7.} $V = 2\times 10^{5}(3e^{-0.5t} - 0.5e^{-3t})/2.5$: $t = 1$:
$1.4\times 10^{5}$; $t = 3$: $5.4\times 10^{4}$; $t = 7$: $7\times
10^{3}$; $t = 14$: $2\times 10^{2}$.
\textbf{8.} $2.4\times 10^{5}e^{-0.5t} = 50$: $t = 2\ln(4800) \approx 17$
hari.
\textbf{9.} Hari 3--10 terletak pada fase lambatnya: kemiringan $\ln V$ di
sana adalah $-\delta$ (dari $5.4\times 10^{4}$ ke $1.6\times 10^{3}$ dalam
tujuh hari, $\delta = 0.50$). Fase cepatnya berlangsung hanya beberapa jam
($1/c$), sehingga $c$ menuntut contoh tiap beberapa jam pada hari pertama,
yang dicocokkan dengan rumus dwieksponensialnya.
\textbf{10.} Sel yang terinfeksi laten, yang melepaskan \omterm{def:b3:virology:virus}{virus} ketika ia
mengaktif kembali; laju peluruhannya $\ln 2/44$ per bulan, yakni sekitar
$5\times 10^{-4}$ per hari.
\textbf{11.} \omterm{def:b3:virology:virus}{Virion} $\ln 2/3 = \qty{0.23}{d}$, sekitar \qty{5.5}{h}; sel
terinfeksi $\ln 2/0.5 = \qty{1.4}{d}$.
\textbf{12.} Muatan yang tetap tampak seperti \omterm{def:b3:virology:virus}{virus} yang tidak berbuat apa-apa;
persamaannya menunjukkannya sebagai keadaan mantap yang di dalamnya $10^{10}$
\omterm{def:b3:virology:virus}{virion} dibuat lalu disingkirkan dan $10^{7}$ sel~T diinfeksi lalu dibunuh
tiap hari.
\textbf{13.} $9700\times 10^{-5} \approx 0.1$ mutasi per genom;
$10^{10}\times 0.1 = 10^{9}$ mutasi sehari.
\textbf{14.} $3\times 9700 \approx \num{29000}$ mutasi tunggal yang mungkin;
masing-masing muncul $10^{9}/\num{29000} \approx 3\times 10^{4}$ kali sehari.
\textbf{15.} Rangkap dua tertentu: $10^{-10}$ per genom, sekitar sekali sehari
tanpa obat; rangkap tiga tertentu: $10^{-15}$, yakni $10^{-5}$ sehari ---
sekali dalam tiga abad.
\textbf{16.} $10^{7}\times 10^{-15} = 10^{-8}$ sehari, $4\times 10^{-6}$
setahun.
\textbf{17.} Dengan satu obat hilang, mutan yang kebal terhadap kedua obat
lain sudah cukup: sebuah rangkap dua tertentu pada $10^{-10}$;
$7\times 10^{7}$ genom dalam seminggu memberi harapan $7\times 10^{-3}$ ---
masih kecil kemungkinannya, kecuali bila muatannya melambung kembali ke arah
$10^{10}$ sehari selama kelalaian itu, ketika ia menjadi satu sehari dan
pelariannya menjadi mungkin. Dosis yang terlewat adalah cara kekebalan datang.
\textbf{18.} Analognya berbagi tapak aktif transkriptase balik,
dan satu mutasi di sana (atau satu himpunan mutasi analog timidin) mengubah
penanganan beberapa di antaranya oleh enzim itu: kekebalan silang. Dua
nukleosida terhitung kurang dari dua obat yang saling bebas, sehingga
regimennya menggabungkan kelas --- sebuah nukleosida, sebuah penghambat
integrase, sebuah \omterm{def:b3:virology:drugs}{penghambat protease} atau sebuah nonnukleosida.
\textbf{19.} $20$ paruh hidup, $880$ bulan, $73$ tahun; menjadi $10^{4}$,
$\log_{2}100 = 6.6$ paruh hidup, $290$ bulan, $24$ tahun.
\textbf{20.} Dari satu \omterm{def:b3:virology:virus}{virion} menjadi $3\times 10^{9}$ pada $R_{0} = 8$ per
generasi: $\ln(3\times 10^{9})/\ln 8 = 10.5$ generasi, yakni tiga minggu.
\textbf{21.} Kejut lalu bunuh: aktifkanlah kembali sel latennya di bawah
terapi supaya mereka mati atau dibunuh --- penghalangnya: tidak ada bahan yang
mengaktifkan semuanya kembali, dan sel yang teraktifkan kembali tidak terbunuh
dengan andal. Penyuntingan gen: potonglah keluar provirusnya atau musnahkanlah
reseptor CCR5 di dalam sel pasiennya --- penghalangnya: menjangkau tiap sel.
(Mengganti sumsumnya dengan sel donor yang tidak ber-CCR5 telah menyembuhkan
segelintir pasien pada risiko yang hanya diterima bila cangkoknya memang
diperlukan.)
\textbf{22.} $1 - 1/R_{0}$: \qty{67}{\%} bagi $R_{0} = 3$, \qty{80}{\%} bagi
$5$.
\textbf{23.} Terapi menetapkan $\beta \to 0$ di dalam pasiennya, muatannya
turun menurut teoremanya hingga hampir nol, dan penularan, yang
sebanding dengan muatannya, berhenti; mengobati tiap orang terinfeksi mendorong
$R_{0}$ populasinya di bawah satu.
\textbf{24.} $f = (1 - 1/R_{0})/E$: dengan $R_{0} = 4$ dan $E = 0.7$, $f =
0.75/0.7 = 1.07$ --- lebih daripada semua orang: \omterm{def:b3:virology:drugs}{vaksinnya} sendiri tidak dapat
mencapai ambangnya; dengan $E = 0.9$, $f = 0.83$.
\textbf{25.} Sekitar $10^{10}$ \omterm{def:b3:virology:virus}{virion} dihasilkan sehari; tak terdeteksi pada
sekitar hari $17$; sekitar $4\times 10^{-6}$ mutan rangkap tiga setahun pada
terapi.
\end{solution}
